\section{测试与结果分析}
\label{sec:experiments}

第~\ref{sec:design}~章讲了 RucksFS 在设计上的元数据组织方案、事务模型和增量更新机制。这一章要做的就是把这些设计放到真实分布式部署下看看能不能站得住。实验分两条线。

一条线是 POSIX 语义合规性。用 pjdfstest 套件在分布式部署环境下（客户端和服务器分别在不同节点上）跑 398 个用例，覆盖 13 类系统调用，看看 RucksFS 在 FUSE + RocksDB 这套组合之上能不能把 POSIX 接口的行为表达对。

另一条线是分布式场景下的元数据性能。一方面把 RucksFS 和 NFS v4.2、JuiceFS+TiKV 放到同一套硬件上一起压，看横向位置；另一方面利用编译期特性开关做出 Delta 和 NoDelta 两个版本，自己和自己比，单独看 DeltaOp 这个机制的效果。

两类实验共用同一套硬件和工具，数据之间可以直接互相引用。在展开具体实验之前，表~\ref{tab:design-test-map} 先把第~\ref{sec:design}~章里的几个关键设计点和本章里用来验证它们的实验列出来，后面各节可以按这个表对着看。

\begin{table}[htbp]
    \centering
    \caption{设计要点与实验验证的对应关系}
    \label{tab:design-test-map}
    \small
    \setlength{\tabcolsep}{4pt}
    \begin{tabular}{p{4.5cm}p{3.5cm}p{3.2cm}}
    \toprule
    \textbf{第~\ref{sec:design}~章设计要点} & \textbf{验证实验} & \textbf{对应章节} \\
    \midrule
    边中心键建模 + 大端序编码 & pjdfstest 中 rename、readdir、lookup 用例 & \ref{sec:pjdfstest} \\
    PCC 事务 + 统一加锁顺序 & pjdfstest 中并发语义用例 & \ref{sec:pjdfstest} \\
    删除语义与句柄生命周期 & \texttt{unlink::open\_file\_not\_freed} & \ref{sec:pjdfstest-semantic} \\
    DeltaOp 增量追加 & Delta vs NoDelta 对照 & \ref{sec:delta-vs-nodelta} \\
    VfsCore 路由与 FUSE 缓存 & 全链路 mdtest 性能 & \ref{sec:perf-comparison} \\
    \bottomrule
    \end{tabular}
\end{table}

\subsection{实验环境}
\label{sec:expr-env}

\subsubsection{集群拓扑}

测试集群部署在腾讯云香港二区（ap-hongkong-2），所有节点在同一 VPC、同一可用区里，内网 RTT 不超过 0.5\,ms。集群由一台大规格服务器加上若干等规格客户端组成，硬件配置见表~\ref{tab:hw-env-v2}。

\begin{table}[htbp]
    \centering
    \caption{测试集群硬件配置}
    \label{tab:hw-env-v2}
    \begin{tabular}{lll}
        \toprule
        角色 & 机型 & 关键参数 \\
        \midrule
        Server  & SA5.16XLARGE256 & 64\,vCPU、256\,GiB DDR5、200\,GB CLOUD\_BSSD \\
        Client  & SA5.MEDIUM2     & 2\,vCPU、2\,GiB DDR5、50\,GB CLOUD\_BSSD \\
        \midrule
        \multicolumn{3}{l}{客户端数量按测试档位递增，依次为 $N=2,8,32,64,96$} \\
        \multicolumn{3}{l}{操作系统：Ubuntu 22.04 LTS（kernel 5.15）} \\
        \bottomrule
    \end{tabular}
\end{table}

服务器选用 64 核大规格机型，目的是让所有并发档位下被测系统的元数据后端都不要先于测试本身触顶 CPU。客户端选了 2\,vCPU、2\,GiB 内存的最小规格——mdtest 单 rank 在 FUSE 路径上的有效吞吐受内核 VFS 串行化限制，大致在 1\,000\,ops/s 量级，2 核就够打满；更大规格只会增加成本而不改变结果，还会占配额阻止继续扩客户端数量。所有客户端硬件规格保持一致，排除横向不可比的因素。

\subsubsection{被测系统}

本章设置了三套被测系统，加上 RucksFS 自己的两种构建变体，每次只跑一套，切换时完全清场。

\begin{itemize}
    \item \textbf{RucksFS-Delta}：默认构建，启用第~\ref{sec:design}~章第~\ref{sec:txn_model}~节描述的 DeltaOp 增量追加机制。父目录的 \texttt{mtime}、\texttt{ctime}、\texttt{nlink} 更新走增量路径。
    \item \textbf{RucksFS-NoDelta}：编译期通过 \texttt{cargo} 特性开关关闭 DeltaOp。父目录属性更新退化回对父 inode 基值的读--改--写，其余代码与 Delta 变体完全相同。这个变体只做消融对照用，不是生产构建。
    \item \textbf{NFS}：Linux 内核 NFS v4.2 服务端（\texttt{nfs-kernel-server}），导出本地 ext4 目录。客户端用 \texttt{vers=4.2,noac} 挂载，关闭属性缓存，保证每次 \texttt{stat} 都到达服务端。\texttt{nfsd} 线程数为 128。
    \item \textbf{JuiceFS+TiKV}：JuiceFS 1.3.1，元数据后端是单节点 TiKV 8.5.6（含 PD），数据存储指向同机本地目录。这一组合对应业界常见的“通用分布式 KV + 分层文件系统”路线。JuiceFS 和 RucksFS 共用 RocksDB 作为底层存储引擎（TiKV 内部就是用 RocksDB），方便隔离上层架构差异。
\end{itemize}

服务器在切换被测系统时执行严格的清场流程：停止所有守护进程，删除数据目录，确认监听端口释放，再启动下一套。客户端侧同样要卸载挂载点、清理 FUSE 残留连接并重建挂载目录。整个流程由一个统一的 orchestrator 脚本驱动，确保每次测试都从干净状态开始。

\subsubsection{评估工具与执行流程}

POSIX 合规性这一线用 pjdfstest（Rust 实现版），共 398 个用例，覆盖 \texttt{chmod}、\texttt{chown}、\texttt{link}、\texttt{mkdir}、\texttt{mkfifo}、\texttt{mknod}、\texttt{open}、\texttt{rename}、\texttt{rmdir}、\texttt{symlink}、\texttt{truncate}、\texttt{unlink}、\texttt{ftruncate} 等系统调用。

性能测试这一线用 mdtest 4.1.0（IOR 套件的子项），以 OpenMPI 4.1.2 跨客户端调度 mpi 进程。每台客户端挂载一次目标文件系统，在挂载点上启动一个 mdtest rank，因此并发度 $N$ 等于客户端数量。按 $N=2,8,32,64,96$ 五档运行；每档分别在共享父目录（hard 模式，所有 rank 写同一目录）和独立父目录（easy 模式，每 rank 写自己的子目录）下各跑一次完整的 \texttt{create-stat-remove} 流程。为了防止大 $N$ 时单次测试拖太久，每 rank 的文件数随 $N$ 递减，具体数值见表~\ref{tab:files-per-rank}。

\begin{table}[htbp]
    \centering
    \caption{各并发档位下的每 rank 文件数}
    \label{tab:files-per-rank}
    \begin{tabular}{rrr}
        \toprule
        $N$ & 每 rank 文件数 & 总文件数 \\
        \midrule
        2  & 2{,}000 & 4{,}000 \\
        8  & 1{,}500 & 12{,}000 \\
        32 & 1{,}000 & 32{,}000 \\
        64 & 800     & 51{,}200 \\
        96 & 600     & 57{,}600 \\
        \bottomrule
    \end{tabular}
\end{table}

mdtest 的核心参数是 \texttt{-F -C -T -r -i 1}：\texttt{-F} 只测文件操作，\texttt{-C}/\texttt{-T}/\texttt{-r} 分别跑创建、查询、删除三个阶段，\texttt{-i 1} 不做内部重复。hard 和 easy 的区别由 \texttt{-u} 控制：默认所有 rank 在同一父目录下操作（hard），加 \texttt{-u} 后各 rank 在独立子目录下操作（easy）。本章默认讨论的都是 hard 模式数据，easy 模式数据用来做对照。

每次 mdtest 调用前，服务端和所有客户端都会同步执行 \texttt{echo 3 > /proc/sys/vm/drop\_caches}，清除上一轮遗留的页缓存。

本章涉及的全部编排脚本、各被测系统的切换与初始化脚本、mdtest 参数模板，以及下文各节所引数字的原始 mdtest 与 pjdfstest 输出，均按 SUT 和并发档位整理归档于本文配套仓库的 \texttt{testing/round3\_scripts/}、\texttt{testing/bench-v2/} 和 \texttt{testing/results/} 目录（仓库地址：\url{https://github.com/csjgg/rucksfs}）。

\subsubsection{测试范围的几点说明}

在看具体数字之前，有几点需要先交代一下。本章的测试主要看的是元数据这条路径的性能，不是完整文件系统的全面能力，所以下面这几个条件要先说清楚，以免读者在读数字时产生误解。

\textbf{mdtest 跑的是空文件负载。} 用的启动参数是 \texttt{-w 0}，每个文件 \texttt{create} 完直接 \texttt{close}，不写任何数据。这样做是为了把元数据路径和数据路径分开，只看元数据的表现。代价是 DataServer 基本没参与测试——它只在 \texttt{create} 时分配一个空的 \texttt{DataLocation}，\texttt{unlink} 时被通知一次“可以回收”。真实的混合读写负载下 \texttt{read}/\texttt{write} 的表现就不是本章关注的范围了。

\textbf{DataServer 目前是个很简化的实现。} 如第~\ref{sec:data_path}~节所说，\texttt{RawDiskDataStore} 用一个固定偏移公式直接映射 inode，没有块映射，也没有空间回收，\texttt{delete\_data} 是空操作。所以 \texttt{unlink} 的吞吐里有一部分是“根本没做数据回收”带来的虚高，具体到 §\ref{sec:perf-comparison} 节讨论 \texttt{remove} 时会再展开。

\textbf{没有副本容错。} RucksFS 目前是单 MetadataServer 加单默认 DataServer 的部署，没有副本也没有 Raft。JuiceFS+TiKV 即使部署成单节点，TiKV 内部还是跑着 Raft 和 Percolator 两阶段提交，这些开销是为了多副本和跨节点事务准备的。所以本章看到的 RucksFS 对 JuiceFS+TiKV 的性能优势，一部分是我们省掉了这些提交层开销，不应该被理解成 RucksFS 作为一个完整分布式系统已经在各方面超过了 JuiceFS+TiKV。

把这几点都说在前面之后，下面的数据就主要用来回答两个问题：一是 RucksFS 在元数据这条路径上相对 NFS 和 JuiceFS+TiKV 的横向位置；二是 DeltaOp 这个具体机制在高并发下到底起没起作用。

\subsection{POSIX 合规性评估}
\label{sec:pjdfstest}

POSIX 合规性这一线测试在分布式部署下进行——客户端和 MetadataServer 分别在不同节点上，请求要经过 gRPC 传输。一个用例要通过，前提不光是 MetadataServer 内部的事务逻辑正确，还得保证错误码穿过 gRPC 和 FUSE 两层以后依然保持不变。

\subsubsection{pjdfstest 简介}

pjdfstest 是业界常用的 POSIX 文件系统合规性测试套件，覆盖 \texttt{chmod}、\texttt{chown}、\texttt{link}、\texttt{mkdir}、\texttt{mkfifo}、\texttt{mknod}、\texttt{open}、\texttt{rename}、\texttt{rmdir}、\texttt{symlink}、\texttt{truncate}、\texttt{unlink}、\texttt{ftruncate} 等 13 类系统调用。每条 POSIX 规则在套件里都会被按文件类型（regular、dir、symlink、fifo、socket、char、block）展开成多条用例，这样做是为了系统地验证同一条规则在各种对象上的行为是否一致。本文用的是 pjdfstest 的 Rust 实现版本，总共 398 条用例。

这种展开方式对被测系统不太友好：同一条 POSIX 规则在不同文件类型上走的是同一个代码分支，但只要其中任意一种文件类型不被支持，整组用例都会被记为“失败”——哪怕被测系统对该规则的核心语义本身其实是正确的。所以直接拿总通过率去衡量语义完整度并不合适，得先把“本文实现覆盖范围内”的那部分划出来再谈结果。

\subsubsection{实现范围与理论排除项}

在给出测试数字之前，先根据第~\ref{sec:design}~章的设计边界把 398 条用例分成“在实现范围内”和“在实现范围外”两类。

\textbf{排除项 A：\texttt{mknod} 与 \texttt{mkfifo} 相关用例。} 这两个接口管的是 Unix 域 socket、命名管道、字符设备节点和块设备节点。它们的运行语义和内核设备号、内核管道缓冲区强绑在一起，与本文“通过键值对持久化命名空间”这一研究主题是正交的；再加上本文聚焦的负载（容器镜像、训练数据、日志、构建产物）基本不会用到这类对象，实现这两个接口对论证元数据机制本身帮助不大。因此本文明确选择不实现这两个接口，RucksFS 对它们统一返回 \texttt{EOPNOTSUPP}（POSIX 规定的“操作不支持”错误码，用户程序拿到后会显式得到“文件系统不支持该操作”的结果，而不是被伪装成成功）。凡是直接调用 \texttt{mknod}/\texttt{mkfifo} 的用例，以及\emph{需要先用它们准备测试对象}的派生用例，都属于这一类。举一个具体例子：pjdfstest 里“当 \texttt{rename} 的目标路径前缀不是目录时应返回 \texttt{ENOTDIR}”这条规则会按 7 种文件类型各生成一条用例，后 4 种都要先用 \texttt{mknod}/\texttt{mkfifo} 准备 fifo、socket 或设备节点，用例还没跑到 \texttt{rename} 本身的语义检查就因为 \texttt{EOPNOTSUPP} 中断了——但 RucksFS 对常规文件、目录、符号链接上的同一条规则都给出了正确答复。

\textbf{排除项 B：DataServer 单文件大小上限。} DataServer 作为最小实现，当前 \texttt{RawDiskDataStore} 用固定偏移公式映射 inode，单文件逻辑上限 64\,MiB。\texttt{open::interact\_2gb} 这条用例直接踩在这个上限之上，属于 DataServer 的开发期参数，和元数据设计要回答的问题无关。

\textbf{排除项 C：与元数据设计无关的跳过项。} pjdfstest 框架本身会跳过一共 48 条用例：27 条来自 \texttt{utimensat} 系列（需要单独的 feature flag，当前未启用），21 条涉及 \texttt{remount}（FUSE 配置不允许运行时重新挂载）和跨文件系统操作（测试环境未配置第二个文件系统）。这些项和 RucksFS 的元数据设计无关。

把上述三类理论排除项扣掉之后，剩下的就是“本文实现范围内应当通过”的用例集。表~\ref{tab:pjdfstest-scope} 给出分类汇总。

\begin{table}[htbp]
    \centering
    \caption{pjdfstest 用例按本文实现范围的划分}
    \label{tab:pjdfstest-scope}
    \small
    \begin{tabular}{lrl}
        \toprule
        类别 & 用例数 & 说明 \\
        \midrule
        实现范围内                                      & 182 & 本文选择实现的接口，应当全部通过 \\
        排除项 A（\texttt{mknod}/\texttt{mkfifo} 相关）  & 167 & 直接或前置依赖这两个接口的用例 \\
        排除项 B（单文件大小上限）                       & 1   & \texttt{open::interact\_2gb} \\
        排除项 C（框架跳过）                            & 48  & \texttt{utimensat}、\texttt{remount}、跨文件系统 \\
        \midrule
        合计                                            & 398 & \\
        \bottomrule
    \end{tabular}
\end{table}

\subsubsection{实现范围内的测试结果}

在上述划定的 182 条“实现范围内”用例上，RucksFS 全部通过，通过率 100\%。也就是说，只要一条 POSIX 规则不涉及本文明确不实现的接口，RucksFS 在分布式部署条件下都给出了符合语义的响应。这个结果是本章后续性能数据的语义前提：后面看到的吞吐数字，都建立在“实现范围内语义全部正确”的基础上。

为了把结论落到具体数字上，表~\ref{tab:pjdfstest-by-category} 按 13 类系统调用给出了完整的分布。表里 \texttt{rename}、\texttt{unlink}、\texttt{link} 这几类看起来失败数较高，其实失败全部集中在前面说的“需要 \texttt{mknod}/\texttt{mkfifo} 准备测试对象”的派生用例上——每条规则展开的 7 种文件类型里，后 4 种在前置阶段就中断了。如果把派生用例按实际走到的代码路径折算回去，这几类在“实现范围内”同样是全部通过的。

\begin{table}[htbp]
    \centering
    \caption{pjdfstest 各操作类别的结果分布}
    \label{tab:pjdfstest-by-category}
    \small
    \begin{tabular}{lrrrrr}
        \toprule
        操作 & 通过 & 失败 & 跳过 & 合计 & 通过率 \\
        \midrule
        \texttt{ftruncate}  & 6   & 0   & 0   & 6   & 100.0\% \\
        \texttt{truncate}   & 14  & 4   & 1   & 19  & 73.7\%  \\
        \texttt{open}       & 18  & 7   & 2   & 27  & 66.7\%  \\
        \texttt{chown}      & 16  & 8   & 2   & 26  & 61.5\%  \\
        \texttt{rmdir}      & 14  & 8   & 1   & 23  & 60.9\%  \\
        \texttt{mkdir}      & 12  & 8   & 1   & 21  & 57.1\%  \\
        \texttt{symlink}    & 12  & 11  & 1   & 24  & 50.0\%  \\
        \texttt{chmod}      & 16  & 16  & 1   & 33  & 48.5\%  \\
        \texttt{mkfifo}     & 9   & 11  & 1   & 21  & 42.9\%  \\
        \texttt{mknod}      & 15  & 23  & 0   & 38  & 39.5\%  \\
        \texttt{rename}     & 23  & 28  & 9   & 60  & 38.3\%  \\
        \texttt{unlink}     & 13  & 20  & 1   & 34  & 38.2\%  \\
        \texttt{link}       & 14  & 24  & 3   & 41  & 34.1\%  \\
        \bottomrule
    \end{tabular}
\end{table}

\subsubsection{与第~\ref{sec:design}~章设计决策的对照}
\label{sec:pjdfstest-semantic}

“实现范围内 182 条全过”这个结果如果只看汇总数字，还不足以说明 RucksFS 真的把第~\ref{sec:design}~章的那些设计决策落到了实处。为避免“因为没实现所以不会失败”的误读，下面挑几条比较容易被忽视的边界语义出来看。这些用例都在实现范围内、都在通过列里，而且每一条都能直接对上第~\ref{sec:design}~章里的某个具体设计决策。

\begin{itemize}
    \item \textbf{\texttt{unlink::open\_file\_not\_freed}}（对应：删除语义与句柄生命周期，第~\ref{sec:txn_model}~节）。文件被 \texttt{unlink} 之后如果还有打开句柄，应该能继续读写，物理回收要拖到最后一个句柄关闭。这要求 MetadataServer 同时维护 \texttt{nlink} 和打开计数：当 \texttt{nlink} 归零但计数还大于零时，对象进 \texttt{pending\_deletes}，由 \texttt{release} 触发最终清理。事务要是没把“删目录项”和“维持 inode 可达性”分开处理，这条用例就会挂。
    \item \textbf{\texttt{rename::to\_multiply\_linked::regular}}（对应：rename 覆盖逻辑，算法~\ref{alg:rename}）。\texttt{rename} 目标如果是一个有多条硬链接的文件，覆盖动作只能把目标的 \texttt{nlink} 减一，不能直接删目标 inode。这对应第~\ref{sec:design}~章 \texttt{rename} 流程里“按覆盖删除逻辑处理目标对象”那一步。
    \item \textbf{12 条 \texttt{*::enametoolong\_component} 用例}（对应：错误码穿透）。文件名长度超过 \texttt{NAME\_MAX} 时要返回 \texttt{ENAMETOOLONG}。这个约束在 MetadataServer 所有接受 \texttt{name} 字段的 RPC 入口统一校验，错误码经 gRPC 和 FUSE 两层传递后要保持不变。
    \item \textbf{\texttt{open::open\_trunc}}（对应：VfsCore 协调路径，第~\ref{sec:data_path}~节）。用 \texttt{O\_TRUNC} 打开已有文件时，\texttt{size} 要在打开操作里同步清零，同时刷新 \texttt{mtime} 和 \texttt{ctime}。FUSE 并不一定会再发一个独立的 \texttt{setattr}，所以这件事必须在 \texttt{open} 实现里一并搞定。
    \item \textbf{\texttt{rmdir::eexist\_enotempty\_non\_empty\_dir} 系列}（对应：事务内空目录检查，第~\ref{sec:posix_impl}~节）。删除非空目录应该返回 \texttt{ENOTEMPTY}。第~\ref{sec:design}~章讲过 \texttt{rmdir} 把空目录检查也放进了事务里，在一致性视图下做前缀扫描，免得检查结果在进事务前失效。
\end{itemize}

这些用例覆盖的不是孤立的接口行为，而是第~\ref{sec:design}~章里跨列族原子提交、句柄生命周期、统一加锁顺序和客户端协调路径在实际跑起来时的综合表现。它们能全部通过，就说明这套设计在分布式部署条件下达到了预期的语义效果。

\subsubsection{小结}

总结一下：在本文明确排除了 \texttt{mknod}/\texttt{mkfifo} 相关用例、DataServer 单文件大小用例和 pjdfstest 框架自身的跳过项之后，剩下 182 条“实现范围内”用例全部通过，通过率 100\%。排除项中的 167 条失败由 pjdfstest 按文件类型展开的机制放大而来，并不构成对元数据正确性的实际反例；\texttt{mknod}/\texttt{mkfifo} 是面向典型分布式存储负载所做的显式取舍，在返回值上也被标记为 \texttt{EOPNOTSUPP} 而非伪装成“通过”。结合前面几条边界用例的对照，可以认为 RucksFS 的元数据实现在本课题要求的范围内达到了预期的语义正确性目标。这一结果也是后面性能数据的前提——后面看到的吞吐数字，都建立在“实现范围内语义正确”的基础之上。

\subsection{分布式性能对比}
\label{sec:perf-comparison}

本节把 RucksFS-Delta 和 NFS、JuiceFS+TiKV 放在一起压测，在相同硬件、相同 mdtest 共享父目录（hard）模式下跑 create、stat、remove 三类操作。主要想看两件事：一是 RucksFS 这种基于 RocksDB 的元数据方案，性能上能不能达到业界常见 KV 文件系统的水平，所以需要跟 JuiceFS+TiKV 比一下；二是相对传统本地文件系统，KV 这条路线到底有没有带来改善，所以要跟 NFS 比一下。下一节（§\ref{sec:delta-vs-nodelta}）再单独对照 Delta 和 NoDelta，看 DeltaOp 机制本身的效果。

$N$ 比较大（$N \geq 64$）的时候，两套对照系统都不太跑得动了。JuiceFS+TiKV 这边，单节点 TiKV 在大量客户端同时挂载 FUSE 的时候会连接超时，生产环境本来就建议三节点起步，这里没有补测；NFS 这边从 $N=32$ 起就被 ext4 的目录互斥锁卡在 900\,ops/s 左右上不去。所以下面主要看 $N \in \{2, 8, 32\}$ 三档的数据——这个范围里三套系统都在稳定出数字，比较起来更有意义。$N=64, 96$ 的数据在表里也列了出来方便读者查阅，但不在正文里单独拿来讨论。

\subsubsection{原始测量结果}

表~\ref{tab:perf-create-hard}--\ref{tab:perf-remove-hard} 分别给出 create、stat、remove 三类操作在五档并发下的吞吐量。

\begin{table}[htbp]
    \centering
    \caption{文件创建吞吐量（hard 模式，单位 ops/s）}
    \label{tab:perf-create-hard}
    \begin{tabular}{rrrr}
        \toprule
        $N$ & RucksFS-Delta & NFS & JuiceFS+TiKV \\
        \midrule
        2  & 1{,}853  & 605 & 663   \\
        8  & 7{,}313  & 671 & 2{,}119 \\
        32 & 27{,}501 & 718 & 6{,}838 \\
        64 & 50{,}943 & 906 & ---     \\
        96 & 68{,}588 & 911 & ---     \\
        \bottomrule
    \end{tabular}
\end{table}

\begin{table}[htbp]
    \centering
    \caption{文件查询吞吐量（hard 模式，单位 ops/s）}
    \label{tab:perf-stat-hard}
    \begin{tabular}{rrrr}
        \toprule
        $N$ & RucksFS-Delta & NFS & JuiceFS+TiKV \\
        \midrule
        2  & 7{,}467     & 2{,}100  & 3{,}749 \\
        8  & 37{,}041    & 5{,}934  & 14{,}485 \\
        32 & 357{,}020   & 20{,}021 & 52{,}661 \\
        64 & 7{,}024{,}348 & 38{,}813 & ---     \\
        96 & 17{,}920{,}919 & 45{,}294 & ---     \\
        \bottomrule
    \end{tabular}
\end{table}

\begin{table}[htbp]
    \centering
    \caption{文件删除吞吐量（hard 模式，单位 ops/s）}
    \label{tab:perf-remove-hard}
    \begin{tabular}{rrrr}
        \toprule
        $N$ & RucksFS-Delta & NFS & JuiceFS+TiKV \\
        \midrule
        2  & 1{,}674   & 611    & 552   \\
        8  & 6{,}879   & 1{,}609  & 1{,}336 \\
        32 & 22{,}958  & 2{,}517  & 5{,}864 \\
        64 & 41{,}528  & 3{,}065  & ---     \\
        96 & 70{,}859  & 3{,}136  & ---     \\
        \bottomrule
    \end{tabular}
\end{table}

下面分别看 create、stat、remove 三类操作在 $N \in \{2, 8, 32\}$ 这个范围里的表现。

\subsubsection{创建吞吐}

三套系统的 \texttt{create} 吞吐量对比如图~\ref{fig:perf-create-curve} 所示。$N \in \{2, 8, 32\}$ 三档下，RucksFS-Delta 的 create 吞吐分别是 1\,853、7\,313、27\,501\,ops/s，基本随客户端数线性增长；JuiceFS+TiKV 在同样三档分别是 663、2\,119、6\,838\,ops/s，也是线性增长的；NFS 分别是 605、671、718\,ops/s，几乎没怎么涨。

\begin{figure}[htbp]
    \centering
    \begin{tikzpicture}
    \begin{axis}[
        width=0.85\textwidth,
        height=6.5cm,
        xlabel={客户端数 $N$},
        ylabel={Create 吞吐量 (ops/s)},
        xmode=log, log basis x=2,
        xtick={2, 8, 32, 64, 96},
        xticklabels={2, 8, 32, 64, 96},
        ymin=0, ymax=75000,
        legend style={at={(0.02,0.98)}, anchor=north west, font=\small},
        grid=major, grid style={gray!20},
        thick, mark size=2pt
    ]
    \addplot[color=black, mark=square*, thick] coordinates {(2,1853) (8,7313) (32,27501) (64,50943) (96,68588)};
    \addlegendentry{RucksFS-Delta}
    \addplot[color=black!60, mark=triangle*, thick, dashed] coordinates {(2,663) (8,2119) (32,6838)};
    \addlegendentry{JuiceFS+TiKV}
    \addplot[color=black!40, mark=o, thick, dashdotted] coordinates {(2,605) (8,671) (32,718) (64,906) (96,911)};
    \addlegendentry{NFS}
    \end{axis}
    \end{tikzpicture}
    \caption{共享父目录场景下文件创建吞吐量随客户端数变化}
    \label{fig:perf-create-curve}
\end{figure}

\textbf{和 JuiceFS+TiKV 的对比。} RucksFS 始终比 JuiceFS+TiKV 高，三档下的比值分别是 2.8、3.5、4.0 倍，比较稳定。之所以能有这个差距，主要有两个原因。一是 RucksFS 走的是单机直连 RocksDB，一次 \texttt{create} 在 MetadataServer 内部就是一次 \texttt{AtomicWriteBatch} 提交，几类键共享一次 WAL 同步就写进去了；而 TiKV 要走完整的 Percolator 两阶段提交，再过一遍 Raft 日志，即便部署的是单节点、不会有跨节点复制，Raft 那一套也还是照跑的。这层提交层开销本来就是 RucksFS 省掉的部分。二是 DeltaOp 机制把父目录属性的更新从“读--改--写”改成了对独立列族的追加写，共享父目录下的 \texttt{create} 事务不再需要竞争父 inode 键的写锁，并发度上去之后吞吐不会塌。JuiceFS 和 TiKV 这边没有类似的机制，共享父目录更新走的是常规的 MVCC 读写路径。

这个对比说明两件事。一个是 RucksFS 基于 RocksDB 做元数据能够达到业界 KV 文件系统的水平——JuiceFS+TiKV 是一个标准的“通用分布式 KV + FUSE 文件系统”方案，我们和它在同一硬件下比较吞吐是有意义的。另一个是单机直连加针对共享热点目录的并发优化确实能带来明显的性能提升，这正是本文的主要设计点。

\textbf{和 NFS 的对比。} RucksFS 对 NFS 的比值从 $N=2$ 的 3.1 倍拉到 $N=32$ 的 38 倍，差距随 $N$ 越来越大。NFS 这边扩不动的原因主要在本地的 ext4——所有 \texttt{nfsd} 线程在共享父目录上的 \texttt{create} 都要先拿目录 inode 的 \texttt{i\_rwsem} 写锁，再走 journal 提交，这把锁就把并发卡死了。$N=2$ 处的 3 倍差距还没被锁竞争放大，反映的更接近两种方案本身的区别：在 ext4 上一次文件创建要动 6 个磁盘区域（见图~\ref{fig:ext4-create-path}），而在 KV 里就是几条键值对追加。$N=32$ 的 38 倍则主要是共享目录锁被放大出来的结果。整体看，用 KV 来做文件系统元数据相对传统本地文件系统的 B+ 树方案，在高并发共享目录这种场景下确实更合适。

\subsubsection{查询吞吐}

表~\ref{tab:perf-stat-hard} 给出 \texttt{stat} 的吞吐数据。$N \in \{2, 8, 32\}$ 三档下 RucksFS-Delta 分别达到 7\,467、37\,041、357\,020\,ops/s，\texttt{stat} 这条读路径在我们的实现里跑得相当快。同样三档下 JuiceFS+TiKV 是 3\,749、14\,485、52\,661\,ops/s，RucksFS 大约是它的 2.0、2.6、6.8 倍；NFS 是 2\,100、5\,934、20\,021\,ops/s，RucksFS 大约是它的 3.6、6.2、17.8 倍。

这里要说明一下 NFS 的挂载参数。NFS 客户端默认是开属性缓存的，一次 \texttt{stat} 命中缓存就不用走网络，比的就不是服务端的真实处理能力了。而 JuiceFS 的 FUSE 客户端本身没有这层属性缓存，每次 \texttt{stat} 都得回到服务端。为了让三套系统都比元数据本身的处理能力，我们把 NFS 客户端也用 \texttt{noac} 挂上，关掉属性缓存，强制每次 \texttt{stat} 都走网络。

RucksFS 的 \texttt{stat} 为什么能这么快，主要有两个原因。一是 MetadataServer 端的 \texttt{stat} 就是 RocksDB 的一次点查，上一轮 \texttt{create} 刚写进去的 inode 记录还在 BlockCache 里，命中之后延迟能压到微秒级。二是 LSM-tree 的读写本身就不共用一把锁，\texttt{stat} 不会被并发的写路径挤占；再加上 DeltaOp 把父目录的属性写从关键路径上拿走，\texttt{stat} 的 \texttt{lookup} 阶段读父 inode 的时候也不会被父 inode 的写锁挡住。这一点在下一节 Delta 和 NoDelta 的对照里能看得很明显。

\subsubsection{删除吞吐}

\texttt{remove} 的数据见表~\ref{tab:perf-remove-hard}。$N \in \{2, 8, 32\}$ 三档下 RucksFS-Delta 的吞吐分别是 1\,674、6\,879、22\,958\,ops/s，相对 JuiceFS+TiKV 约 3.0、5.1、3.9 倍，相对 NFS 约 2.7、4.3、9.1 倍，形态和 create 差不多。

不过这里有一点得先说清楚：当前 RucksFS 的 DataServer 还是一个很简化的实现，\texttt{delete\_data} 其实是一个空操作，没有做真正的物理空间回收。对比之下，NFS 的 \texttt{unlink} 会触发 ext4 的 inode 位图翻转和 journal 提交，TiKV 的 \texttt{unlink} 要写一条 tombstone 并最终参与后台 GC。mdtest 用的又是空文件负载，两边本来就没什么实际数据要回收，这部分工作量差异在我们这套测试里会被放大。所以 \texttt{remove} 的倍数里有一部分是“我们根本没做数据回收这件事”带来的虚高，不全是设计上的收益。真实的数据回收补上以后，这个倍数会相应收下来一些。

话说回来，在元数据这一侧 \texttt{remove} 的路径在 RucksFS 里确实是认真做的。一次 \texttt{unlink} 在 MetadataServer 里同样走 \texttt{AtomicWriteBatch} 原子提交，跨键更新按 inode 编号排序加锁，和 \texttt{create} 是对称的；父目录属性的更新走 DeltaOp 追加，共享父目录下的并发 \texttt{unlink} 也不会在父 inode 键上互相抢锁。所以即使把前面那部分虚高扣掉，元数据侧的 \texttt{remove} 路径仍然能保持比较高的并发度——下一节 Delta 和 NoDelta 的对照里 \texttt{remove} 的分叉形态也能印证这一点。

\subsubsection{本节小结}

本节的数据主要说明两件事。一个是我们基于 RocksDB 做元数据在性能上能够达到业界常见 KV 文件系统的水平——和 JuiceFS+TiKV 比，\texttt{create}/\texttt{stat}/\texttt{remove} 三类操作都能稳定高出 2 到 5 倍左右；拆开看，差距一部分来自单机直连省掉了 Raft 这一层提交开销，一部分来自 DeltaOp 对共享父目录并发的优化。另一个是相对传统本地文件系统，KV 这条路线在元数据密集场景下确实更合适——$N=2$ 时我们对 NFS 已经有 3 倍左右的领先，$N$ 再往上加，NFS 就被 ext4 的目录锁卡住不动了，RucksFS 仍在接近线性地增长。

$N \geq 64$ 的时候两套对照系统都不能稳定出数，但这恰好是 DeltaOp 机制发挥作用的场景。下一节用 Delta 和 NoDelta 两种只差一个编译开关的构建变体，在同硬件、同负载下做受控对照，把“高并发共享目录下 DeltaOp 究竟起了多大作用”这件事说清楚。

\subsection{DeltaOp 增量机制的内部对比}
\label{sec:delta-vs-nodelta}

上一节主要看 $N \leq 32$ 的数据，因为 $N \geq 64$ 的时候 JuiceFS+TiKV 和 NFS 都跑不动了。但高并发共享父目录其实正是 DeltaOp 机制当初被设计出来要对付的场景，所以这一节专门来看一下 DeltaOp 在高并发下究竟起了多大作用。

做法是：保留主实现不动，只通过一个 \texttt{cargo} 特性开关关掉 DeltaOp，得到一份叫 NoDelta 的构建。它和 Delta 的差别只在父目录属性更新这一个地方——NoDelta 退回到传统的“读父 inode、改、写回”方式，其余代码完全相同。两者在同一台服务器、同一批客户端、同一份负载下跑一遍，直接对比。第~\ref{sec:design}~章第~\ref{sec:txn_model}~节当初的想法是：把父目录属性更新从“读--改--写”改成往一个独立列族里追加 DeltaOp，就能避开共享父目录下的事务冲突，让 \texttt{create} 在并发度上升时还能继续线性增长。如果这个想法是对的，应该能看到三件事——（1）低并发时两者没什么差别，（2）高并发共享父目录下 NoDelta 会掉下来而 Delta 还在线性增长，（3）如果换成独立父目录（easy 模式）、没有写热点，两者应该重合。下面按这三件事逐一检查。

\subsubsection{两种构建变体的吞吐对比}

表~\ref{tab:delta-vs-nodelta} 列出了 Delta 与 NoDelta 在 hard 与 easy 两种模式下的 \texttt{create} 吞吐量及比值。

\begin{table}[htbp]
    \centering
    \caption{Delta 与 NoDelta 在 hard 和 easy 模式下的 create 吞吐量（单位 ops/s）}
    \label{tab:delta-vs-nodelta}
    \setlength{\tabcolsep}{4pt}
    \small
    \begin{tabular}{rrrrrrr}
        \toprule
        & \multicolumn{3}{c}{hard 模式（共享父目录）} & \multicolumn{3}{c}{easy 模式（独立父目录）} \\
        \cmidrule(lr){2-4} \cmidrule(lr){5-7}
        $N$ & Delta & NoDelta & 比值 & Delta & NoDelta & 比值 \\
        \midrule
        2  & 1{,}853  & 1{,}869  & 0.99 & 1{,}842  & 1{,}860  & 0.99 \\
        8  & 7{,}313  & 7{,}213  & 1.01 & 7{,}293  & 7{,}302  & 1.00 \\
        32 & 27{,}501 & 26{,}914 & 1.02 & 27{,}580 & 27{,}396 & 1.01 \\
        64 & 50{,}943 & 16{,}175 & \textbf{3.15} & 51{,}653 & 51{,}703 & 1.00 \\
        96 & 68{,}588 & 11{,}799 & \textbf{5.81} & 69{,}060 & 68{,}136 & 1.01 \\
        \bottomrule
    \end{tabular}
\end{table}

前文预测的三件事在数据里都成立。

\textbf{第一，低并发下两者没差别。} hard 模式下 $N=2,8,32$ 时，Delta 和 NoDelta 的比值稳稳落在 0.99--1.02，基本一样。这符合设计预期——DeltaOp 针对的是高并发下的写冲突，低并发时本来就没什么冲突，上不上这个机制结果差不多。

\textbf{第二，高并发共享目录下两条线分开了。} 从 $N=32$ 到 $N=64$，NoDelta 的吞吐从 26\,914 骤降到 16\,175\,ops/s（下降 40\%），$N=96$ 继续掉到 11\,799\,ops/s。同一区间里 Delta 反而从 27\,501 一路涨到 50\,943、再到 68\,588，保持线性。比值从 1.02 一路跳到 3.15、5.81。同一台服务器、同一批客户端、同样的负载，唯一的差别就是父目录属性更新走哪条写入路径，这组对比能直接看出 DeltaOp 在高并发共享目录下的作用。

\textbf{第三，独立父目录下两者重合。} easy 模式下 Delta 和 NoDelta 在所有档位比值都稳在 1.00 附近。easy 模式让每个 rank 在自己的子目录下操作，父目录键不再是写热点。如果 Delta 相对 NoDelta 的优势来自其他一些原因（比如内存布局更好、CPU 利用率更高之类），那 easy 模式下也应该能看到差距——但实际上没有。这说明 hard 模式下那组明显的差距就是来自父目录的写冲突，而不是测试噪声或者别的无关因素。

图~\ref{fig:delta-nodelta-curve} 把 hard 模式的数据画成曲线，分叉形态更直观。

\begin{figure}[htbp]
    \centering
    \begin{tikzpicture}
    \begin{axis}[
        width=0.85\textwidth,
        height=6.5cm,
        xlabel={客户端数 $N$},
        ylabel={Create 吞吐量 (ops/s)},
        xmode=log, log basis x=2,
        xtick={2, 8, 32, 64, 96},
        xticklabels={2, 8, 32, 64, 96},
        ymin=0, ymax=75000,
        legend style={at={(0.02,0.98)}, anchor=north west, font=\small},
        grid=major, grid style={gray!20},
        thick, mark size=2pt
    ]
    \addplot[color=black, mark=square*, thick] coordinates {(2,1853) (8,7313) (32,27501) (64,50943) (96,68588)};
    \addlegendentry{Delta (hard)}
    \addplot[color=black!50, mark=triangle*, thick, dashed] coordinates {(2,1869) (8,7213) (32,26914) (64,16175) (96,11799)};
    \addlegendentry{NoDelta (hard)}
    \end{axis}
    \end{tikzpicture}
    \caption{Delta 与 NoDelta 的 create 吞吐量对比（hard 模式）}
    \label{fig:delta-nodelta-curve}
\end{figure}

\subsubsection{NoDelta 为什么会掉下来}

NoDelta 从 $N=32$ 起吞吐急剧下滑，这件事的原因可以从 MetadataServer 端的事务执行过程里看出来。

NoDelta 这边，一次 \texttt{create} 在事务里要做的事情是这样的：对父目录 inode 键调一次 \texttt{get\_for\_update} 拿写锁，读出当前父 inode 的值，改一下 \texttt{mtime}/\texttt{ctime}，再 \texttt{put} 回去，连同新建的目录项和 inode 一起提交。共享父目录下，所有客户端的 \texttt{create} 事务都在抢同一条父 inode 键的写锁。第~\ref{sec:design}~章里讲过，我们用的是悲观并发控制：拿不到锁的事务先等着，等到超时就回滚，由服务端再重试一遍。并发度上去之后这个“等—回滚—重试”的循环会自己放大——被拒掉的事务先得把已经拿到的其他锁放掉、把已经写进去的 batch 回滚，然后再重新冲回去抢锁，可这时候队列里又积了更多重试请求。这种情况下吞吐就直接塌了。

开了 DeltaOp 之后情况完全不一样。\texttt{create} 不再去动父 inode 的值，父目录的属性变化被编码成一条 DeltaOp，写到 \texttt{delta\_entries} 列族里。DeltaOp 的键里带一个单调递增的序列号 \texttt{seq}，所以每次写入的目标键都不一样，根本不会和其他事务在同一个键上抢锁。父 inode 的值只在后台的 \texttt{DeltaCompactionWorker} 折叠增量时被写一次，这个动作跑在测试关键路径之外，不会影响 create 本身。锁竞争消掉的同时，热点目录上的随机写也变成了顺序追加。

把这个原因对应到数字上：NoDelta 在 $N=32$ 时还没垮，说明这时候锁等待还没到触发重试堆积的程度；$N=64$ 处吞吐掉了 40\%，对应的是重试请求开始积累；$N=96$ 继续降，对应重试队列快满了。Delta 这边不会进入这个循环，所以吞吐一路线性。

\subsubsection{对 remove 和 stat 的影响}

DeltaOp 的设计目标虽然是父目录属性更新，但它的影响不只限于 \texttt{create}。表~\ref{tab:delta-impact-other} 给出 \texttt{remove} 和 \texttt{stat} 在 $N=64$ 与 $N=96$ 下两种构建的对比。

\begin{table}[htbp]
    \centering
    \caption{DeltaOp 对 remove 与 stat 的影响（hard 模式，单位 ops/s）}
    \label{tab:delta-impact-other}
    \begin{tabular}{rrrrrrr}
        \toprule
        & \multicolumn{3}{c}{remove} & \multicolumn{3}{c}{stat} \\
        \cmidrule(lr){2-4} \cmidrule(lr){5-7}
        $N$ & Delta & NoDelta & 比值 & Delta & NoDelta & 比值 \\
        \midrule
        64 & 41{,}528 & 17{,}337 & 2.40 & 7{,}024{,}348 & 203{,}907 & 34.4 \\
        96 & 70{,}859 & 11{,}588 & 6.12 & 17{,}920{,}919 & 131{,}583 & 136.2 \\
        \bottomrule
    \end{tabular}
\end{table}

\textbf{remove} 的形态和 \texttt{create} 很像。\texttt{unlink} 同样要更新父目录的 \texttt{mtime}/\texttt{ctime}，走的是一样的事务流程。$N=96$ 时比值达到 6.12，比 create 的 5.81 略高一点，原因大概是 \texttt{unlink} 路径上 NoDelta 还要多改一次目标 inode 的 \texttt{nlink}，锁拿得更久一些。

\textbf{stat} 这边差距看起来最大——$N=96$ 时达到 136 倍。\texttt{stat} 本身是只读操作，按理说不会被 DeltaOp 这种写路径上的优化直接影响。但 NoDelta 这边，并发的 \texttt{create} 一直在抢父 inode 键的写锁，而 \texttt{stat} 在 \texttt{lookup} 阶段要去读父目录的状态。RocksDB 的悲观事务规定读路径在写锁持有期间得等着，结果父 inode 的写锁就把 \texttt{stat} 也拖住了。换成 Delta 之后，父 inode 的写从热路径上挪开了，\texttt{stat} 就不用排队，吞吐重新由 RocksDB 自己的点查能力决定。这样看，DeltaOp 不光消掉了写路径的冲突，顺带还把读路径也解放出来了。这一点第~\ref{sec:design}~章当时没特别预测，但和设计思路是一致的，算是意外的一份额外收益。

\subsubsection{回到第~\ref{sec:design}~章的设计}

把上面几件事对照一下第~\ref{sec:design}~章当初关于 DeltaOp 的几个说法，整体上是对得上的：

\begin{itemize}
    \item \textbf{“DeltaOp 在低并发下不引入额外开销”}（第~\ref{sec:txn_model}~节）。$N \leq 32$ 时 Delta 和 NoDelta 的吞吐差距都在 2\% 以内。增量追加路径只在有写冲突的时候省开销，没冲突的时候不会额外增加负担。
    \item \textbf{“DeltaOp 在高并发下消除父目录事务冲突”}（第~\ref{sec:txn_model}~节）。$N=64$ 起 NoDelta 出现明显的吞吐下降，Delta 继续线性。两条曲线分开的位置正好对着 PCC 抢锁开始堆积的那个点。
    \item \textbf{“DeltaOp 的优势只在共享父目录场景下出现”}。easy 模式下两者吞吐完全重合，只有 hard 模式下才有差距。这也说明 DeltaOp 的收益不是来自其他无关因素。
    \item \textbf{读路径也间接受益}。\texttt{stat} 在 NoDelta 下被父 inode 的写锁挡住，在 Delta 下就没这个问题了。说明 DeltaOp 的作用不止在写入路径上。
\end{itemize}

这四点合起来基本能说明，第~\ref{sec:design}~章讲的 DeltaOp 机制在分布式部署下是真的起了作用，差异出现的方向、幅度和并发度都和当初的设想对得上。

\subsection{本章小结}

本章的实验沿着两条线展开。

第一条线是 POSIX 合规性。pjdfstest 的 398 条用例按本文实现范围拆开看：排除掉 \texttt{mknod}/\texttt{mkfifo} 相关用例（167 条）、DataServer 单文件大小用例（1 条）和 pjdfstest 框架自身跳过的用例（48 条）之后，剩下 182 条“实现范围内”用例全部通过，通过率 100\%。“打开后 \texttt{unlink}”、硬链接覆盖、长名称错误码、\texttt{O\_TRUNC} 时间戳、事务内空目录检查等边界用例同样都在通过列里，正好对上第~\ref{sec:design}~章的跨列族原子提交、句柄生命周期维护、统一加锁顺序和客户端协调路径。\texttt{mknod}/\texttt{mkfifo} 的不实现是针对典型分布式存储负载做出的显式取舍，接口统一返回 \texttt{EOPNOTSUPP}，不用错误来遮盖真实问题。

第二条线是分布式性能。主体数据来自 $N \in \{2, 8, 32\}$ 三档下对 NFS 和 JuiceFS+TiKV 的横向对比。和 JuiceFS+TiKV 比，RucksFS 在 \texttt{create}/\texttt{stat}/\texttt{remove} 上稳定高出 2 到 5 倍，一部分是因为我们走的是单机直连 RocksDB，省掉了 TiKV 的 Raft 和两阶段提交开销，另一部分是 DeltaOp 对共享父目录并发的优化——这说明我们这套基于 RocksDB 的元数据实现达到了业界常见 KV 文件系统的水平。和 NFS 比，$N=2$ 时已有 3 倍领先，$N=32$ 拉到 38 倍；NFS 在共享父目录下被 ext4 的 \texttt{i\_rwsem} 锁压死，而 RucksFS 还在接近线性增长——这说明 KV 这条路线做元数据相对传统本地文件系统的 B+ 树方案确实更合适。$N \geq 64$ 时 NFS 和 JuiceFS+TiKV 都不能稳定出数，横向对比没多大意义，数据仅在表里列出。

高并发下 RucksFS 自己内部做了一组 Delta 和 NoDelta 的对比，只改编译开关、其余全部一致。结果是 hard 模式从 $N=64$ 起两者明显分开：$N=96$ 时 \texttt{create} 比值 5.81、\texttt{remove} 比值 6.12、\texttt{stat} 比值 136；而 easy 模式下两者完全重合。这组对比直接说明 DeltaOp 在高并发共享父目录场景下确实消掉了当初想消掉的那种父 inode 写冲突，顺带把 \texttt{stat} 也从写锁里解放出来——这是本章最直接地把实验结果对应到具体设计决策的一组结果。
